\section{密度响应函数与介电函数}

\subsection{密度--密度响应}

定义
\begin{equation}
  \chi(\mathbf{q},\omega)
  \equiv\chi_{nn}(\mathbf{q},\omega)
  =-\frac{\mathrm{i}}{\hbar}
  \int_0^{\infty}\mathrm{d}t\,
  e^{\mathrm{i}\omega t}
  \bigl\langle
  \bigl[\hat n_{\mathbf{q}}(t),\hat n_{-\mathbf{q}}(0)\bigr]
  \bigr\rangle.
\label{eq:chinndef}
\end{equation}
若外加标势使
\begin{equation}
  H'(t)
  =\frac{1}{V}\sum_{\mathbf{q}}
  \hat n_{-\mathbf{q}}\,
  v_{\mathrm{ext}}(\mathbf{q},t),
\end{equation}
则诱导密度
\begin{equation}
  \delta n(\mathbf{q},\omega)
  =\chi(\mathbf{q},\omega)\,
  v_{\mathrm{ext}}(\mathbf{q},\omega).
\label{eq:dn_chi}
\end{equation}
（不同书对 $v_{\mathrm{ext}}$ 是否含电荷因子 $-e$ 的约定不同；以下与 $v_{\mathbf{q}}=4\pi e^2/q^2$ 一致地讨论电子气。）

\subsection{Proper / improper 极化与介电函数}

电子气中总势 = 外势 + Hartree 诱导势。定义 \textbf{不可约}（proper）密度响应 $\chi_{\mathrm{irr}}$ 为对\textit{总}势的响应：
\begin{equation}
  \delta n=\chi_{\mathrm{irr}}\,v_{\mathrm{tot}},
  \qquad
  v_{\mathrm{tot}}=v_{\mathrm{ext}}+v_{\mathbf{q}}\delta n.
\end{equation}
与 \eqref{eq:dn_chi} 比较得
\begin{equation}
  \chi
  =\frac{\chi_{\mathrm{irr}}}
  {1-v_{\mathbf{q}}\chi_{\mathrm{irr}}}.
\label{eq:chi_irr}
\end{equation}
介电函数定义为
\begin{equation}
  \varepsilon(\mathbf{q},\omega)
  =\frac{v_{\mathrm{ext}}}{v_{\mathrm{tot}}}
  =1-v_{\mathbf{q}}\chi_{\mathrm{irr}}(\mathbf{q},\omega).
\label{eq:eps_def}
\end{equation}
从而屏蔽相互作用
\begin{equation}
  W(\mathbf{q},\omega)
  =\frac{v_{\mathbf{q}}}{\varepsilon(\mathbf{q},\omega)}
  =\frac{v_{\mathbf{q}}}{1-v_{\mathbf{q}}\chi_{\mathrm{irr}}}.
\end{equation}
集体模式（等离激元）满足 $\varepsilon(\mathbf{q},\omega)=0$。

压缩率求和规则必须作用于\textbf{不可约}响应的 $q\to0$ 极限（或等价地作用于 $\varepsilon(q,0)$ 的适当组合）：因 $v_q$ 发散，全响应 $\chi_{nn}(q\to0,0)\to0$，不能直接等于 $-n^2K$。电容测量读取的是经 $v_q$ 屏蔽后的可观测量（GV §5.2.3）。

RPA（$\chi_{\mathrm{irr}}=\chi_0$）与局域场因子 $G(\mathbf{q},\omega)$ 的展开见本章后两节。